% Transcription — fonds Alexandre Grothendieck, shelfmark 105, batch 1 (pages 1-20).
% Established from the facsimile archives/batches/105/batch-01.pdf (a mirror of
% https://grothendieck.umontpellier.fr/105.pdf), per the skill
% transcribe-grothendieck. The batch file carries no cover sheet: PDF page N
% is page N of the folder (the Montpellier footer prints « 1 » ... « 20 »),
% and the archivists' pencil numbers bottom-left (bottom-right on page 6)
% agree wherever legible (« 1 », « 2 », « 3 », « 4 », « 6 », « 9 », « 11 »,
% « 13 », « 15 », « 17 », « 19 »).
%
% Pass: Opus 5.5 (claude-opus-5-5), 2026-09-23 — first pass, unchecked against the pages by a human.
%
% Ten of the twenty pages are transcribed: 1, 2, 4, 6, 9, 11, 13, 15, 17, 19.
% Absent: page 3, a sheet of pink card carrying only a pencilled name struck
% through and a few ink spots, nothing mathematical; pages 5, 7, 10, 12, 14,
% 16, 18 and 20, sheets of computer listing (job logs, tables, a banner page)
% carrying nothing of his — page 18 has two pencilled words at its head,
% not mathematics; page 8, a blank sheet. The gaps in the numbering are the
% record. Several of the transcribed leaves are themselves written over
% listing paper (its print shows through); nothing is taken from the print.
%
% Two loose scratch leaves, in French, open the batch (pages 1 and 2); they
% carry no pagination of his. The run proper, in English, begins on page 4
% and is paginated by him at the head of each leaf: page 4 is his 1, 6 his 2,
% 9 his 3, 11 his 4, 13 his 5, 15 his 6, 17 his 7, 19 his 8. Batch 2 opens
% (page 21) on his 9, so the run continues across the boundary unbroken.
% The inventory title « [Champs (stacks) 2] » is the archivists'; both
% \section headings here are ours (bracketed).
%
% What the batch is.
%   1      Scratch (French): generating maps ∅ → e and ∂ × S^n(N) → I × S^n(N)
%          plus « cônes » over them; a Lemma: RLP with respect to them makes
%          X W_∞-aspherical (a), resp. X → Y W_∞-proper / W_∞-smooth (b);
%          a heavily overwritten corollary. Much of the prose is lost.
%   2      Scratch (French), cancelled by two long diagonals: Ar(C) → C × C
%          should have fibration-type properties; it is neither a
%          categorical fibration nor cofibration; wish that p be a
%          W_∞-fibration; Question: do W_∞-fibrations and the cofibrations
%          defined by LLP form a Quillen structure on (Cat)?
%   4-9    (his 1-3) A a topos, W ⊂ Fl(A) weakly saturated (a, b, c);
%          (TC)_0 ⊂ C_0 given; TF, F by RLP, then C, TC by LLP; the pairs
%          TC ↔ F, C ↔ TF; stabilities; TC = C ∩ W~, TF = F ∩ W~ with
%          W~ = { pi | p ∈ TF, i ∈ TC }; for i: A → B, f: X → Y the map
%          α(i,f): Hom(B,X) → Hom(i,f) = Hom(B,Y) ×_{Hom(A,Y)} Hom(A,X);
%          equivalence of (i) C, TC stable by ? × T and (ii) TF, F stable by
%          Hom(T,?), with proof; (iii) on α(i,f); an NB axiomatising this
%          through Φ^*, Φ_* and a pair ⊗, Hom tied by the Cartan relation;
%          the box product i ⊠ i'.
%   11     (his 4) the formula α(i', α(i,f)) ≃ α(i' ⊠ i, f); a table of four
%          equivalences between ⊠-stabilities and α-conditions; Example:
%          ⊗ = cartesian product, C = mono, (TC)_0 the monomorphisms
%          (B_0 × e) ∪ (A_0 × I) ⊂ B_0 × I for I in a bunch of W-aspherical
%          intervals with a section e.
%   13     (his 5) the square computing i' ⊠ h(i_0, I, e) = h(i' ⊠ i_0, I, e);
%          « We do not know yet M1 ... nor M5 »; programme: identify W~,
%          prove W~ = W, the dependence of (TC)_0 on the « bunch » of I's.
%   15-17  (his 6-7) W~ ⊂ W reduced to TC ⊂ W, TF ⊂ W; axiom ① (L_A → e)
%          ∈ UW; ② (I → e) ∈ UW for all I in the bunch; ③ stability of
%          monomorphisms in W under cobase change; ④ stability of W under
%          filtering direct limits; ⑤ W stable by direct factors; retract
%          argument; « Thus ① to ⑤ imply W~ ⊂ W ».
%   19     (his 8) the converse W ⊂ W~: reduced to F ∩ W ⊂ TF or C ∩ W ⊂ TC;
%          RLP w.r.t. (TC)_0 versus all monomorphisms in W; two cases wanted:
%          a) W = W_A for a (local?) W-test category A, b) A = A_0 × T.
%
% Notation fixed for the batch. Fl(A) as \mathrm{Fl}(A); the underlined Hom
% (internal) as \underline{\mathrm{Hom}}, the plain Hom as \mathrm{Hom};
% his circled cross for the product « ⊗ » as \otimes, his boxed cross for
% the pushout-product of two maps as \boxtimes; W with a tilde as
% \widetilde{W}; UW as \mathrm{U}W; his cursive capital Φ (drawn close to a
% script L, as in batch 2) as \Phi; the object « T » of ? × T and Hom(T,?),
% drawn as a T whose stem is crossed, as T (on page 19 T is a test category,
% a different T); I (interval) and e (final object, and the section of I)
% as he writes them; e_A the final object of A; the basic localizer
% « W » underlined as \underline{W}; circled numbers ① ... ⑤ as he circles
% them. Words he underlines in prose are set \emph. His « p.r. »,
% « compl. », « asph. », « s.th. », « i.e. » are kept.
%
% The dating is the inventory's own for the folder, copied verbatim; the
% title in \foldertitle is the archivists' (square brackets), not his.
\documentclass[11pt,a4paper]{article}
% Le préambule commun à toutes les transcriptions du fonds Grothendieck.
%
% Il tient en une idée : **l'incertitude se compose, elle ne se résout pas.**
% Une transcription qui lisse un mot douteux a détruit la seule chose pour
% laquelle on la faisait — un lecteur doit pouvoir savoir, à chaque endroit,
% ce qui était sur le feuillet et ce qui a été ajouté. Les six macros qui
% suivent sont donc l'appareil critique, et non de la décoration.
%
% Les mêmes commandes sont reconnues par `scripts/render.mjs`, qui produit la
% vue de lecture du site. Une macro ajoutée ici doit l'être aussi là-bas,
% faute de quoi le rendu échouera bruyamment — ce qui est voulu : un rendu
% silencieusement faux est pire qu'un rendu absent.

\ProvidesPackage{grothendieck}[2026/08/08 Transcriptions du fonds Grothendieck]

\usepackage{amsmath,amssymb,amsthm}
\usepackage{mathrsfs}
\usepackage{stmaryrd}
% Les diagrammes commutatifs sont partout dans ce fonds : c'est la langue
% même de la géométrie algébrique de Grothendieck, et une transcription qui
% les rendrait en prose ne transcrirait rien.
\usepackage{tikz-cd}
% Français par défaut — c'est la langue des feuillets — mais l'anglais est
% chargé aussi : la traduction et le résumé sont des documents anglais, et la
% typographie française (espaces avant les deux-points) y serait une faute.
% Ces documents basculent par \selectlanguage{english} en tête de corps.
\usepackage[english,french]{babel}
\usepackage{fontspec}
\usepackage{geometry}
\geometry{a4paper,margin=25mm,marginparwidth=28mm}
\usepackage{marginnote}
\usepackage{xcolor}
% ulem, not soul. `soul`'s \st and \ul are built on a character-by-character
% scan that XeLaTeX's font machinery breaks: the document fails with an
% undefined control sequence rather than degrading. ulem does the same two jobs
% and is Unicode-engine safe. `normalem` stops it hijacking \emph, which the
% transcriptions use for Grothendieck's own emphasis.
\usepackage[normalem]{ulem}
\usepackage{hyperref}
\hypersetup{colorlinks=true,linkcolor=black,urlcolor=black,pdfborder={0 0 0}}

% --- Métadonnées du lot -----------------------------------------------------
% Reprises telles quelles par le script de rendu, qui les affiche en tête de
% la vue de lecture. Elles fixent ce que le fichier prétend couvrir : sans
% elles, une transcription flotte sans point d'attache dans un fonds de seize
% mille feuillets.

\newcommand{\folder}[1]{\gdef\grothFolder{#1}}
\newcommand{\batch}[1]{\gdef\grothBatch{#1}}
\newcommand{\leaves}[2]{\gdef\grothFirst{#1}\gdef\grothLast{#2}}
\newcommand{\foldertitle}[1]{\gdef\grothFoldertitle{#1}}
\gdef\grothFolder{?}\gdef\grothBatch{?}\gdef\grothFirst{?}\gdef\grothLast{?}\gdef\grothFoldertitle{}

% --- Le résumé --------------------------------------------------------------
%
% Il ouvre la lecture modernisée, sous le titre « Résumé », et il est composé
% en petit. Ce n'est pas un ornement : c'est le seul passage du document qui
% ne soit pas des mathématiques, et un lecteur doit pouvoir le reconnaître
% avant de l'avoir lu — puis le sauter s'il connaît déjà le sujet, ou s'y
% arrêter s'il ne le connaît pas. Le filet qui le suit marque la reprise du
% corps.

\newenvironment{resume}
  {\small\setlength{\parskip}{0.45em}}
  {\par\medskip\hrule\medskip}

% --- Mots-clés ---------------------------------------------------------------
%
% En anglais, en clôture du résumé de la lecture modernisée : le vocabulaire
% moderne sous lequel chercher ce que les feuillets construisent. Le manifeste
% du site (`npm run manifest`) les extrait pour étiqueter la cote entière —
% cette ligne est la source unique des tags, il n'y a pas de fichier à côté.
\newcommand{\keywords}[1]{\par\smallskip\noindent
  {\footnotesize\textsc{Keywords} --- \textit{#1}}\par}

% --- L'appareil critique ----------------------------------------------------

% Le numéro de feuillet, en marge. C'est l'ancre qui permet de régler un
% désaccord sur une formule en regardant *un* feuillet plutôt qu'une chemise
% de six cents. Le site s'en sert aussi pour tourner les pages du fac-similé
% quand on fait défiler la transcription.
\newcommand{\leaf}[1]{%
  \marginnote{\footnotesize\color{gray!70}\textsf{#1}}%
  \label{leaf:#1}\ignorespaces}

% Illisible : on ne devine pas. Un mot inventé qui se lit comme les autres est
% le pire résultat possible.
\newcommand{\ill}{\textcolor{red!60!black}{[\,\ldots\,]}}

% Lecture douteuse : le mot est donné, et signalé comme incertain.
\newcommand{\uncertain}[1]{\uline{#1}}

% Ajout de l'éditeur — une lettre manquante, un mot sous-entendu.
\newcommand{\add}[1]{\textcolor{blue!50!black}{[#1]}}

% Biffé par Grothendieck lui-même. Ce qu'il a rayé fait partie du document :
% une rature dit souvent où la pensée a changé de direction.
\newcommand{\struck}[1]{\sout{#1}}

% Note du transcripteur, sur l'état matériel du feuillet.
\newcommand{\note}[1]{\par\smallskip{\small\color{gray!60!black}\textit{[#1]}}\par\smallskip}

% Note marginale de Grothendieck, distincte du corps du texte.
\newcommand{\marginal}[1]{\par\smallskip\noindent\hspace{1em}%
  {\small\color{gray!60!black}#1}\par\smallskip}

% --- Titre ------------------------------------------------------------------

\AtBeginDocument{%
  \begin{center}
    {\large\bfseries Fonds Alexandre Grothendieck --- cote n\textsuperscript{o}~\grothFolder}\\[2pt]
    {\itshape\grothFoldertitle}\\[4pt]
    {\small feuillets \grothFirst--\grothLast\ (lot \grothBatch)}\\[2pt]
    {\footnotesize\color{gray!60!black}Transcription établie d'après le fac-similé numérisé
     par l'Université de Montpellier.\\
     Les lectures incertaines sont soulignées, les illisibles notées [\,\ldots\,],
     les ajouts entre crochets.}
  \end{center}
  \vspace{1em}}

\endinput


\folder{105}
\batch{1}
\pages{1}{20}
\dating{[à partir de 1982]}
\watermark{Édition de démonstration}
\foldertitle{[Champs (stacks) 2] : notes manuscrites (s.d.).}

\begin{document}

\section{[Deux brouillons : un lemme de relèvement, la flèche $\mathrm{Ar}(C) \to C \times C$]}
\note{titre de l'éditeur, entre crochets : les deux feuillets n'en portent
aucun et ne sont pas paginés de sa main ; ils sont en français, à la
différence de la suite.}

\page{1}
\note{en haut à gauche, un signe en forme de Z barré suivi de « a », non
identifié}

\[
  \partial \xrightarrow{\ i\ } I
\]
\note{sous $I$, une courte mention illisible entre parenthèses ; plus bas,
deux lignes biffées :}
\struck{\ill{}\,$\Pi_0(\partial) \geq 2$} \quad \struck{$\varnothing$ \ill{}}

\[
  \left\{
  \begin{array}{l}
    \varnothing \to e \\[2pt]
    \partial \times S^n(N) \longrightarrow I \times S^n(N)
  \end{array}
  \right.
\]
$+$ cônes \ill{} (compl. droits) sur les précédents (donc aussi
$\{0\} \to \Delta^1$ resp. $\{1\} \to \Delta^1$)

\[
  \underline{\mathrm{Hom}}(T, X)
\]
\note{la lettre, très appuyée, est lue $T$ ; elle pourrait aussi être le $I$
du haut de la page}

\emph{Lemme} a) Si $X \to e$ a la propriété RLP p.r. à $\varnothing \to e$
\emph{et} les $\partial \times S^n(N) \to I \times S^n(N)$ pour
$N \in \mathcal{N}$ (\struck{\ill{}} \uncertain{partie finie} de
$\mathbb{N}$), alors $X$ est $W_\infty$-asph\supplied{érique}.

b) Si $X \to Y$ a la RLP p.r. aux cônes \ill{} (compl. droits) p.r. aux
morphismes précédents, alors $X \to Y$ est $W_\infty$-propre
(resp. $W_\infty$-\uncertain{lisse}).

\emph{Cor.} a) Si $X \to Y$ a la propriété RLP p.r. aux
\uncertain{(compl. gauches)} \ill{}, \struck{\ill{}} alors c'est \ill{}
morphisme $W_\infty$-propre \ill{} à $W_\infty$-fibres asphériques
(compl. \uncertain{lisse}).
\note{le $Y$ de $X \to Y$ est écrit au-dessus d'un $e$ ; le passage est
surchargé et en partie couvert par une grosse tache d'encre, et plusieurs
mots biffés s'y superposent (« \uncertain{remarques} », « \uncertain{droites} »)}

b) Si $f : X \to Y$ a la propriété RLP p.r. \struck{\ill{}} \ill{}
\uncertain{(donc aussi)} \ill{} qui sont des $W_\infty$, alors $f$ est
$W_\infty$-propre (resp. $W_\infty$-\uncertain{lisse}). S'il a \ill{}
propriétés \ill{} pour les $W_\infty$ \ill{} \struck{\ill{}} \ill{}
cat. \uncertain{finies} \ill{} $W_\infty$-\uncertain{parfait}.
\note{les deux dernières lignes, tachées, ne se lisent que par fragments}

\page{2}
\note{feuillet déchiré, écrit dans la largeur, barré en entier par deux
longs traits obliques qui se rejoignent en haut ; il se lit sous la rature}

\uncertain{Notamment} on voudrait que
\[
  \mathrm{Ar}(C) \xrightarrow{\ p\ } C \times C
\]
\note{« Ar » est souligné}
jouisse des propriétés de type fibration \struck{\ill{}}. Or
1) au sens catégorique, ce n'est ni une fibration, ni une cofibration, et
\ill{} les \struck{propriétés} \uncertain{composées}
$\underline{s} \circ p$, $\underline{b} \circ p : \mathrm{Ar}(C) \to C$
ne le sont pas.

2) De \ill{}, \struck{\ill{}} $p$ n'est ni \uncertain{coh.} propre, ni
\uncertain{coh.} lisse. Peut-être au moins \struck{\ill{}}
$\underline{s} \circ p$ et $\underline{b} \circ p$ \struck{\ill{}}
sont-ils \struck{\ill{}} l'un ou l'autre ? Mais ce sont plutôt les
prop\supplied{riétés} de $p$ lui-même qui importent !

3) Je souhaiterais que $p$ soit une \struck{\ill{}}\,$W_\infty$-fibration
(« \uncertain{fibrations cohomologiques} »). C'est une chose à vérifier.

\emph{Question} Les $W_\infty$-fibrations, et les notions correspondantes de
cofibrations (par la L.L.P), \struck{\ill{}} font-ils une
\struck{des} structure de Quillen sur $(\mathrm{Cat})$ ?

\section{[Paires de relèvement dans un topos, la classe $\widetilde{W}$ et sa comparaison avec $W$]}
\note{titre de l'éditeur, entre crochets. À partir d'ici le texte est en
anglais, tel qu'il l'écrit, et paginé de sa main 1 à 8 ; la suite (p.\ 9
de l'auteur) ouvre le lot 2, page 21.}

\page{4}
\note{p.\ 1 de l'auteur.}
$A$ a topos

$W \subset \mathrm{Fl}(A)$ set of arrows, \add{weakly} saturated
\note{« weakly » est ajouté au-dessus de la ligne}
i.e.
\[
  \left\{
  \begin{array}{l}
    \text{a) isom. are in } W \\
    \text{b) if } f = vu \text{, and two of } u, v, f \text{ are in } W \text{, so is the third} \\
    \text{c) if } X \underset{g}{\overset{f}{\rightleftarrows}} Y \text{ with } gf = \mathrm{id}_X,\ fg \in W \text{, then } f, g \in W
  \end{array}
  \right.
\]
\note{b) se lit « if $f = \uncertain{vu}$, and two \ill{} of $u, v, f$ in
$W$, so is the third » ; la fin de la ligne mord sur le bord du feuillet}
plus other conditions we are going to list, as we need them.

\struck{\ill{}} Let moreover \struck{$C_0(TC)_0$}
\[
  (TC)_0 \subset C_0 \ (\subset \mathrm{Fl}(A))
\]
(with no extra assumptions for the time being). We define \struck{$TF_0$}
\[
  TF \subset F \ (\subset \mathrm{Fl}(A))
\]
by the right lifting property (RLP) with respect to \struck{$(TC)_0$}
$C_0$ and $(TC)_0$ respectively, and then
\[
  \begin{array}{ccc}
    TC & \subset & C \quad (\subset \mathrm{Fl}(A)) \\
    \cup & & \cup \\
    (TC)_0 & \subset & C_0
  \end{array}
\]
by the \struck{right} left lifting property (LLP) with respect to $F$ and
$TF$ respectively. We have pairs

\begin{tikzcd}
  TC \arrow[r, leftrightarrow] \arrow[d, hook] & F \\
  C \arrow[r, leftrightarrow] & TF \arrow[u, hook]
\end{tikzcd}

\note{les deux inclusions verticales sont tracées comme des crochets
d'inclusion ; un trait vertical sépare les deux colonnes}

in which each \add{one} determines the other by RLP or LLP.
\struck{Moreover} Moreover
\begin{itemize}
  \item[a)] $F$, $TF$ stable by composition, base change, direct factors
  \item[b)] $C$, $TC$ stable by composition, cobase change, direct factors
\end{itemize}

\marginal{NB $C$ should be \uncertain{called} « cofibrations », $TC$ for
« trivial cofibrations », $F$ for « fibrations », and $TF$ for « trivial
fibrations »}
\note{note écrite en oblique dans la marge de droite, à hauteur de la
définition de $(TC)_0$}

\page{6}
\note{p.\ 2 de l'auteur.}
\struck{c) $X \to Y$ is $\in W$}
\[
  TC = C \cap \widetilde{W} \qquad TF = F \cap \widetilde{W}
\]
where
\[
  \widetilde{W} \subset \mathrm{Fl}(A), \qquad
  \widetilde{W} = \bigl\{\, f = pi \bigm| p \in TF,\ i \in TC \,\bigr\}
\]
\note{le second $\widetilde{W}$ est indiqué par un signe d'égalité vertical
sous le premier}

\emph{NB} These facts do only make use of the existence of \struck{fin.}
amalgamated sums and fibered products in $A$ (not of the fact that $A$ is
a topos, — and for the time being, $W$ has not been used). We used only the
symmetric relationship between \struck{\ill{}} $(TC, F)$, and $(C, TF)$,
\uncertain{and plus} $TC \subset C$ ($\Leftrightarrow TF \subset F$), not the
way they were obtained by some \add{given} $(TC)_0 \subset C_0$.

d) \struck{Assume $TC$, if $i \in C$, $f \in$} consider \struck{but}, for
any
\[
  i : A \to B \qquad f : X \to Y
\]
\note{sous $i$ et $f$, biffés : « $i \in C$ », « $f \in F$ »}
\struck{consider}
\[
  \underline{\mathrm{Hom}}(B, X) \xrightarrow{\ \alpha(i,f)\ }
  \underline{\mathrm{Hom}}(i, f) \overset{\mathrm{def}}{=}
  \underline{\mathrm{Hom}}(B, Y)
  \times_{\underline{\mathrm{Hom}}(A, Y)}
  \underline{\mathrm{Hom}}(A, X)
\]

\struck{Then $\alpha(i,f) \in$} Then we have the following
\struck{equivalences}
\begin{itemize}
  \item[(i)] $C$, $TC$ stable by cartesian products $? \times T$
  \item[(ii)] $TF$, \struck{$T$}$F$ stable by $\underline{\mathrm{Hom}}(T, ?)$
  \item[(iii)] If $i \in C$, $f \in F$, then $\alpha(i,f) \in F$, and if
  moreover $i \in TC$ or $f \in TF$, then $\alpha(i,f) \in TF$.
\end{itemize}
\note{une accolade réunit (i) et (ii), et une flèche double part de (iii)
vers elle ; en marge gauche, à hauteur de (iii) :}
\marginal{under suitable assumptions}

\page{9}
\note{p.\ 3 de l'auteur.}
\emph{Proof} (i) $\Rightarrow$ (ii) Assume (i), let $f : X \to Y$,
$f \in F$, prove \struck{$\underline{\mathrm{Hom}}(T, f)$}
\[
  \bigl( \underline{\mathrm{Hom}}(T, f) : \underline{\mathrm{Hom}}(T, X)
  \to \underline{\mathrm{Hom}}(T, Y) \bigr) \in F
\]
i.e.\ \struck{LLP} RLP with respect to any $i : A \to B$ \add{in $TC$} —
this amounts to RLP of $f$ with respect to
$i \times \mathrm{id}_T : A \times T \to B \times T$, OK as this is
\add{still} in $TC$. Same for $f \in TF$, using stability of $C$ under
$\times T$.

(ii) $\Rightarrow$ (i) Assume (ii), let $i : A \to B$, \struck{given}
$\in C$, prove $i \times T : A \times T \to B \times T$ is still in $C$,
i.e.\ has \struck{right} LLP with respect to any $f \in TF$, but this
amounts to $i$ having the same with respect to
$\underline{\mathrm{Hom}}(T, f)$, which is still in $TF$ by assumption OK.

\struck{(i)+(ii) $\Rightarrow$ (iii)} \emph{NB} The equivalence
(i) $\Leftrightarrow$ (ii) could be further axiomatized as a relation
between $\Phi^*$, $\Phi_*$ associated by LLP and RLP, when we have two
functors $\otimes$, $\underline{\mathrm{Hom}}$ tied by the Cartan relation
\[
  \mathrm{Hom}(A \otimes T, X) \simeq
  \mathrm{Hom}(A, \underline{\mathrm{Hom}}(T, X))
\]

\struck{(i)+(ii) $\Rightarrow$ (iii)} Moreover, when we have stability of
$\Phi^*$ \struck{under $? \otimes T$} (when $\Phi^*$ is either $C$ or
$TC$) or equivalently of $\Phi_*$ under $\underline{\mathrm{Hom}}(T, ?)$,
then for $i : A \to B$ in $\Phi^*$, $f : X \to Y$ in $\Phi$,
\uncertain{(or in $\Phi_*$ ?)} we have that \struck{$\underline{\mathrm{Hom}}(i,f)$}
\[
  \alpha(i,f) :
  \underline{\mathrm{Hom}}(B, X) \to \underline{\mathrm{Hom}}(i, f)
  \quad \text{is in } \emph{TF}
\]
\note{« is in $TF$ » est souligné, et une parenthèse ajoutée le rattache à
la suite}
Indeed, This means it has RLP with respect to any $j : A' \to B'$ in $C$
\add{$[\Phi^*]$}, i.e.
\[
  \mathrm{Hom}(B', \underline{\mathrm{Hom}}(B, X)) \longrightarrow
  \mathrm{Hom}(j, \alpha(i,f))
\]
\note{devant $\alpha(i,f)$, un « $\underline{\mathrm{Hom}}(i,f)$ » biffé}
is surjection. \struck{But the first \ill{}} can be viewed as \struck{$\mathrm{Hom}(B' \otimes B, X)$}
\[
  \mathrm{Hom}(B' \otimes B, X) \longrightarrow
  \mathrm{Hom}(j \boxtimes i, f)
\]
\note{au-dessus de « $j \boxtimes i$ », deux essais biffés de la même
notation}
where \struck{$i \boxtimes i : A' \otimes A \to B' \otimes B$}
\[
  j \boxtimes i : \ A' \otimes B \amalg_{A' \otimes A} B' \otimes A
  \longrightarrow B' \otimes B
\]
\struck{and} Thus we have to check that if $i, i' \in \Phi^*$, then
$i \boxtimes i'$ is in $\Phi^*$ [OK if $\Phi^*$ are monomorphisms, say, and
$\otimes$ = cartesian product]
\marginal{check \ill{} \uncertain{compatibility} \ill{} — and conversely}

Coming back to case of $(C, TF)$, $(TC, F)$, this shows that if $i \in C$,
$f \in TF$, then \struck{\ill{}} $\alpha(i,f)$ is in $TF$, \add{and} if
$i$ is in $TC$ and $f$ in $F$, then \struck{\ill{}} $\alpha(i,f)$ is in
$F$,

\page{11}
\note{p.\ 4 de l'auteur.}
provided we check that if $i, i' \in TC$ \marginal{(mono in $C$ ?)}, then
$i' \boxtimes i$ is in $TC$ \marginal{(mono in $C$ ?)}.
\note{les deux remarques entre parenthèses sont écrites au-dessus de la
ligne et rattachées au texte par des accolades}
\struck{But} If this is so even if we assume only \emph{one} of the two
arrows $i, i'$ in $TC$, the other being in $C$, then we get even that
\[
  i \in TC,\ f \in F \Rightarrow \alpha(i,f) \in TF
  \quad\text{and}\quad
  i \in C,\ f \in TF \Rightarrow \alpha(i,f) \in TF .
\]
We just need formally
\[
  \alpha(i', \alpha(i,f)) \simeq \alpha(i' \boxtimes i, f)
\]
\note{le $\alpha$ intérieur est souligné}

\emph{To sum up} we get equivalences
\[
  \begin{array}{rcl}
    C \otimes T \subset C & \Longleftrightarrow & TF \text{ stable by } \underline{\mathrm{Hom}}(T, ?) \\
    TC \otimes T \subset TC & \Longleftrightarrow & F \text{ stable by } \underline{\mathrm{Hom}}(T, ?) \\
    C \boxtimes C \subset C & \Longleftrightarrow & \bigl( i \in C,\ f \in TF \Rightarrow \alpha(i,f) \in TF \bigr) \\
    \ \boxtimes TC \subset TC & \Longleftrightarrow & \bigl( i \in TC,\ f \in F \Rightarrow \alpha(i,f) \in TF \bigr) \\
    TC \boxtimes C \subset TC & \Longleftrightarrow & \bigl( i \in C,\ f \in F \Rightarrow \alpha(i,f) \in F \bigr)
  \end{array}
\]
\note{les membres de gauche des deux premières lignes sont écrits
par-dessus « \struck{stable by $? \otimes T$} », biffé ; la troisième est
écrite sous « \struck{$C$ stable by $i \boxtimes i'$} » ; le membre de gauche
de la quatrième est noyé sous des hachures, au-dessous desquelles on lit
« $TC \boxtimes TC$ » ou \uncertain{« $C \boxtimes TC$ »} ; dans son membre
de droite, $TF$ est écrit sur une autre lettre. En face de la troisième,
dans la marge, un mot biffé : \struck{\ill{}}}

\emph{Example} \struck{$A$ a topos} $A$ a topos, $\otimes$ = cartesian
product, $\underline{\mathrm{Hom}}$ the usual $\underline{\mathrm{Hom}}$
\add{(\uncertain{inner} \ill{})},
$C = \text{mono}$, $(TC)_0$ = monomorphisms of the type
\struck{\ill{}}
\[
  (B_0 \times e) \cup (A_0 \times I)
  \overset{i}{\subset} B_0 \times I ,
\]
where $A_0 \xrightarrow{\ i_0\ } B_0$ any mono, $I$ any
\emph{$W$-aspheric} \uncertain{object} over $e_A$, endowed with
\uncertain{a} section $e$.
\note{« $(B_0 \times e)$ » est ajouté au-dessus d'un début biffé ; en
face, dans la marge de droite :}
\marginal{given bunch of $I$'s with given sections}

$C$ stable by $? \times T$ : trivial, hence \uncertain{$F$} stable by
$\underline{\mathrm{Hom}}(T, ?)$
\note{on attend ici « $TF$ » ; la lettre qui précède $F$, s'il y en a une,
est à peine marquée}

$(TC)_0$ stable by $? \times T$ (replace $i_0$ by
$i_0 \times \mathrm{id}_T$) hence $F$ stable by
$\underline{\mathrm{Hom}}(T, ?)$, hence $TC$ stable by $? \times T$.

\struck{$C$ stable by $i' \boxtimes i$ : trivial}
$C \boxtimes C \subset C$ : trivial

To check $TC \boxtimes C \subset TC$, $C \boxtimes TC \subset TC$.
Consider $i$ defined by $i_0 : A_0 \to B_0$, $I$, $e$ (i.e.\
$i = h(i_0, \uncertain{I, e})$) and $i' : A' \to B'$, consider the
diagrams

\begin{tikzcd}
  A' \times A \arrow[r, hook] \arrow[d, hook] & B' \times A \arrow[d, hook] \\
  A' \times B \arrow[r, hook] & B' \times B = B' \times B_0 \times I
\end{tikzcd}

\note{plusieurs lettres du carré sont surchargées ($A'$ sur $A$, $B$ sur un
autre signe)}

\page{13}
\note{p.\ 5 de l'auteur.}

\begin{tikzcd}[column sep=small, row sep=small, nodes={font=\scriptsize}]
  ((A' \times B_0) \times e) \cup (A' \times A_0 \times I) \arrow[r] \arrow[d, "{h(i_0 \times \mathrm{id}_{A'},\, I,\, e)}"] & ((B' \times B_0) \times e) \cup ((B' \times A_0) \times I) \arrow[d, "{h(i_0 \times \mathrm{id}_{B'},\, I,\, e)}"] \\
  (A' \times B_0) \times I \arrow[r, hook] & B' \times B_0 \times I
\end{tikzcd}

\note{les deux étiquettes $h(\ldots)$ sont écrites entre les flèches
verticales, chacune reliée à la sienne par une accolade ; sous le coin
inférieur gauche, un trait ondulé ; au-dessous, deux débuts biffés :
\struck{$A' \times B_0$}, \struck{$A' \boxtimes h(i$}}

On trouve
\[
  i' \boxtimes \bigl( i = h(i_0, I, e) \bigr) = h(i' \boxtimes i_0, I, e),
\]
\emph{\uncertain{vérifier}} \ill{} et la formule symétrique, donc
\uncertain{OK}.
\note{ces deux lignes sont en français}

\emph{We do not know yet} M1 (factorization of any $f$ as $pi$, with either
$i \in C$, $p \in TF$, or $i \in TC$, $p \in F$), nor M5 (saturation of
$\widetilde{W}$ — indeed we need only check a), b) to have a closed model
category). The factorization $f = pi$ is the usual argument for constructing
injectives, using reasonable properties of existence of \add{small}
generating set of $A$, and of filtering direct limits. The lifting property
may again be axiomatized in terms of $\Phi^*$, $\Phi_*$, in a topos or a
more general category — I'll not for the moment work out the general
formulation, and \struck{to \ill{}} I'll assume the factorizability
\struck{\ill{}} in the two ways.
\marginal{to work out}

We now want to identify $\widetilde{W}$, and check conditions b) of
saturation. We want even prove $\widetilde{W} = W$ — now it becomes
essential to \struck{\ill{}} make clear how the definition of $(TC)_0$
\emph{depends} on $W$, taking as « bunch » of \struck{intervals} all
intervals $I$ which are $W$-aspheric over $e$. But this may make
difficulties for factorization,

\page{15}
\note{p.\ 6 de l'auteur.}
because this bunch is not « small », we'll therefore take \emph{some}
smaller bunch, and see what we need in it to be able to conclude
$\widetilde{W} = W$.

$\widetilde{W} \overset{?}{\subset} W$. We are reduced to proving
$TC \subset W$, $TF \subset W$. As $TF$ is defined in terms of
$C$ = monomorphisms, the inclusion
\[
  \boxed{TF \subset W}
\]
\uncertain{may} be taken as an axiom on $W$. But in case we got an
\add{separating} interval $I$ which is $TF$ over $e$ i.e.\ « injective
object », \struck{we may take} this inclusion is equivalent with
\struck{$(I \to e) \in \mathrm{U}W$ i.e. $(I \times X \to X)$ is in $W$ for any $X$}
\note{la formule est encadrée, et l'encadré barré d'un trait ondulé en haut et en bas}

Thus, the axiom is \uncertain{really}
\[
  \text{①}\qquad \boxed{(L_A \to e) \in \mathrm{U}W}
\]
\note{au-dessus de $L$, une flèche et deux mots : \uncertain{« Lawvere
element »}}
\uncertain{or equivalently} \struck{$\exists$ separating interval $I$} with
$(I \to e) \in$ \ill{}
\note{variante ajoutée à droite de l'encadré, en partie biffée}
This implies that any $L$-homotopism is in $W$. But for \struck{any} any
separating interval $I$, it is clear that
\[
  f \in TF \Longrightarrow f \text{ is a } I\text{-homotopism}.
\]
Next we got to prove $TC \subset W$. We need first of all
\[
  (TC)_0 \subset W .
\]
\struck{\ill{}}

\begin{tikzcd}[column sep=small, row sep=small, nodes={font=\scriptsize}]
  A_0 \times \{e\} \arrow[r, hook] \arrow[d, "W"'] & B_0 \times \{e\} \arrow[d, hook, "f" description] \arrow[dr, "W"] & \\
  A_0 \times I \arrow[r, hook] & A \arrow[r, hook] & B = B_0 \times I
\end{tikzcd}

\note{la flèche verticale médiane porte « $f$ » et « $W$ » ; le premier
terme, $A_0 \times \{e\}$, est précédé d'une parenthèse biffée}

We need for this
\[
  \left|
  \begin{array}{l}
    \text{if } S \overset{i}{\hookrightarrow} T \text{ mono, and} \\
    S \overset{g}{\hookrightarrow} S' \text{ mono} \in W \\
    \text{then in } \quad
    \begin{array}{ccc}
      S & \overset{i}{\hookrightarrow} & T \\
      {\scriptstyle g} \downarrow & & \downarrow {\scriptstyle f} \\
      S' & \overset{i'}{\hookrightarrow} & S' \amalg_S T = T'
    \end{array} \\
    f \in W
  \end{array}
  \right.
\]
\note{au pied du feuillet, « and », qui annonce la suite}

\marginal{\struck{strict} \uncertain{strongly} \ill{} topos : if this is
\ill{}, \ill{} merely $\{0, 1\}$ can be \ill{} in $A$ — but this is always
\uncertain{so}, because the \uncertain{lower} object is injective !!!}
\note{note écrite en oblique dans la marge de gauche, reliée par une flèche
à « Thus, the axiom is really » ; lue par fragments seulement}

\page{17}
\note{p.\ 7 de l'auteur.}
\struck{An fact} that for any \struck{interval} $I$ \struck{\ill{}} in our
bunch,
\[
  A \times e \to A \times I \quad \text{is in } W \text{ for any } A,
\]
\struck{i.e. $(e \to I)$}.
This is \add{equivalent to} \struck{\ill{}} assuming
$A \times I \to A$ is in $W$ (using saturation) i.e.
\[
  \text{②}\qquad
  \boxed{(I \to e) \in \mathrm{U}W. \text{ for any } I \text{ in the bunch}}
\]

We have to prove that by saturating $(TC)_0$ into $TC$, we still remain
within $W$. But when we use factorization of
\[
  f : A \overset{TC}{\hookrightarrow} B \qquad f \in TC
\]
as
\[
  f = pi : \ A \underset{TC}{\hookrightarrow} B' \underset{TF}{\longrightarrow} B ,
\]
the map $A \overset{i}{\hookrightarrow} B'$ is not only in $TC$, but
\add{also} in $W$ if we assume a property of stability of $W$ under
\add{filtering direct} limits, \marginal{④}
and moreover \struck{under} of \struck{stability of monomorphisms}
stability of monomorphisms in $W$ by cobase extensions \marginal{③} (which
is anyhow necessary for $W = \widetilde{W}$). \struck{We get}
\note{le groupe « stability of monomorphisms in $W$ by cobase extensions »
est encadré ; en marge gauche, en face : « \ill{} \uncertain{explicit} »}

Thus any map $f$ factors into $pi$, with $p \in F$,
$i \in TC \cap W$, and as $f \in TC$, to conclude $f \in W$, we only
have to prove $p \in W$. But we get the lifting property for $(f, p)$ a
section $s$ of $B' \to B$, such that $sf = i$, thus $f$ is a retract of
(= direct factor) of $i \in W$,

\begin{tikzcd}
  & B' \arrow[d, "p"] \\
  A \arrow[ur, "i"] \arrow[r, hook, "f"'] & B \arrow[u, bend right=40, "s"']
\end{tikzcd}

\begin{tikzcd}
  A \arrow[r, hook, "i"] \arrow[d, "="'] & B' \arrow[d, "p"] \\
  A \arrow[r, hook, "f"'] & B \arrow[u, bend right=40, "s"']
\end{tikzcd}

\note{deux figures en marge gauche : dans la première, la flèche $i$ porte
« $W$ » et « (and $TC$) », $p$ porte « $F$ », $f$ porte « $TC$ » ; dans la
seconde, $i$ porte « $W$ », $p$ porte « $F$ », et l'identité de gauche est
dessinée comme deux flèches opposées entre deux signes « = »}

Thus we need \uncertain{④}

⑤ $W$ stable by direct factors (not a consequence of
\uncertain{saturation} ?) as a necessary condition for
$W = \widetilde{W}$ (because this will \ill{} we have a closed model
structure with $W$, \struck{\ill{}} and $\widetilde{W}$ is stable by
retractions). Thus ① to ⑤ imply $\widetilde{W} \subset W$.
\note{le quatrième chiffre entouré, après « Thus we need », est mal formé ;
lu \uncertain{④}}

\page{19}
\note{p.\ 8 de l'auteur.}
We still have to see if $W \subset \widetilde{W}$. For the
$\widetilde{W} \subset W$ inclusion \add{(in terms of $W$)} we had to make
restrictions on the $I$'s (condition ②), now to get $\widetilde{W}$ large
enough, i.e.\ $TC$ large enough, we must make an assumption which will
ensure that the bunch of $I$'s is large enough — a natural condition would
be that it contains a separating interval. We know already there are
« eligible » intervals, namely $W$-aspheric over $e$ and
\add{separating}, such as $L$.

Let $f \in W$, consider factorization
\[
  X \overset{i}{\longrightarrow} X'
  \overset{p}{\longrightarrow} Y
  \qquad i \in TC,\ p \in F \ \ (\text{or } i \in C,\ p \in TF)
\]
\note{sous la flèche $i$, un « $TC$ » biffé}
we know already that $i \in W$ (resp! $p \in W$) therefore $p \in W$
(resp.\ $i \in W$) — hence we are reduced to \struck{Then one of} proving
either one of the two inclusions
\[
  \begin{array}{ll}
    F \cap W \subset TF & (\text{necessary for } W \subset \widetilde{W}, \text{ as } TF = F \cap \widetilde{W}) \\
    C \cap W \subset TC & (\text{id})
  \end{array}
\]
which are \struck{\ill{}} equivalent. \struck{\ill{}} \ill{} to the
condition \add{[first ?]}
\[
  W \subset \widetilde{W} \quad (\text{or equivalently, }
  \widetilde{W} \cap W = W)
\]
\add{[\ill{}]}.
\note{devant le second $\widetilde{W}$, un premier essai biffé ; au-dessous, une insertion entre crochets, illisible}
The second means that if $f : X \to Y$ (has the \struck{right} RLP with
respect to \struck{\ill{}} $(TC)_0$, i.e.\ all $h(i_0, I, e)$), then it has
the RLP with respect to \emph{all} monomorphisms which are in $W$.
[ ? ??? ]

We would like this to hold at least in the following two cases :
\begin{itemize}
  \item[a)] $W = \underline{W}_A$ ($\underline{W}$ a basic localizer),
  $A$ a (\uncertain{local} ?) $\underline{W}$-test category,
  \item[b)] $A = A_0 \times T$, with $T$ a test category,
  $\widehat{A} \simeq \underline{\mathrm{Hom}}(A_0^{\mathrm{op}}, \widehat{T})
  \simeq (A_0 \times T)^{\wedge}$, $W = \underline{W}_{A/A_0}$
  ($\underline{W}$ a basic localizer), $\underline{W}_{A/A_0}$ is defined
  componentwise via $\underline{W}_T$.
\end{itemize}
\note{dans b), le groupe « $\widehat{A} \simeq \ldots$ » est ajouté au-dessus
de la ligne ; en marge gauche, en bas : « \ill{} \uncertain{test} \ill{}
$A = A^{\ill{}}$ »}

\marginal{but this cannot possibly be enough, as it would \ill{}
\struck{\ill{}} just by itself, given just one particular \ill{}, hence
one particular $F$, hence one particular \ill{} on $W$ ...}
\note{note écrite en oblique dans la marge de gauche, à hauteur de « a
natural condition would be that it contains a separating interval »}

\end{document}
